\begin{filecontents*}{\jobname.xmpdata}
\Title{The Ledger of Meluhha --- Indus Valley Script as Metrological Accounting Code}
\Author{Rajeshkumar Venugopal}
\Language{en-US}
\Subject{The Indus Valley script is a metrological cargo-tag system encoding commodity class, weight tier, and merchant identity. Cross-corpus decode of Tamil Nadu potsherds through IVC seal codebook confirms jar goods and Mesopotamia route from real CISI data.}
\Keywords{Indus Valley, Harappan script, metrological accounting, Bronze Age trade, Meluhha, weight systems, sign frequency, cargo tag, CISI, Alloy, SQLite}
\Publisher{Third Buyer Advisory}
\Date{2026-04-16}
\end{filecontents*}

\RequirePackage{luatex85}
\documentclass[12pt]{article}

\usepackage{fontspec}
\setmainfont{TeX Gyre Pagella}

\usepackage{amsmath,amssymb,amsthm}
\usepackage[margin=1.2in,top=1.3in,bottom=1.3in]{geometry}
\usepackage{microtype}
\usepackage{setspace}
\usepackage{xspace}
\usepackage{graphicx}

\usepackage[bottom,hang,flushmargin]{footmisc}
\usepackage{booktabs}
\usepackage{array}
\usepackage{longtable}
\usepackage{adjustbox}
\usepackage{siunitx}

\usepackage{enumitem}
\usepackage[font=small,labelfont=bf,skip=6pt]{caption}
\usepackage{titlesec}
\usepackage{fancyhdr}
\usepackage{xcolor}
\usepackage{luacode}

\usepackage{pgfplots}
\pgfplotsset{compat=1.18}
\usepackage{tikz}
\usetikzlibrary{patterns,positioning,arrows.meta,calc}

\PassOptionsToPackage{
  colorlinks=true,
  linkcolor=black,
  citecolor=black,
  urlcolor=black
}{hyperref}
\usepackage{pdfx}
\usepackage{cleveref}
\usepackage{orcidlink}

\directlua{dofile("citations.lua")}
\newcommand{\smartcite}[1]{\directlua{_smartcite("#1")}}

\definecolor{TBAred}{HTML}{b5121b}

\setlength{\headheight}{13.6pt}
\emergencystretch=4em
\setstretch{1.15}
\setlength{\parskip}{0.4em}
\setlength{\parindent}{1.5em}
\setlist{noitemsep,topsep=4pt}
\renewcommand{\arraystretch}{1.2}

\newcommand{\sign}[1]{\textsc{S-#1}}
\newcommand{\mfield}[1]{\texttt{\small #1}}
\newcommand{\shorttitle}{The Ledger of Meluhha}

\titleformat{\section}{\large\bfseries}{}{0em}{}[\color{TBAred}\titlerule]
\titleformat{\subsection}{\normalsize\bfseries}{}{0em}{}
\titleformat{\paragraph}[runin]{\normalsize\bfseries}{}{0em}{}[.]

\pagestyle{fancy}
\fancyhf{}
\renewcommand{\headrulewidth}{0.3pt}
\fancyhead[L]{\small\href{https://thirdbuyeradvisory.org/}{thirdbuyeradvisory.org}}
\fancyhead[R]{\small\textit{\shorttitle}}
\fancyfoot[C]{\small\thepage}

\title{\textbf{The Ledger of Meluhha}\\
  \large\textit{Indus Valley Script as Metrological Accounting Code}}

\author{Rajeshkumar Venugopal \orcidlink{0009-0002-1838-5976}\\
  \small Third Buyer Advisory LLC, Waterford, Michigan\\
  \small\href{mailto:vrajeshkumar@gmail.com}{vrajeshkumar@gmail.com}}

\date{April 2026\\[4pt]\small Peer review requested.}

\begin{document}

\maketitle
\setlength{\skip\footins}{10pt}

\noindent\textbf{Abstract.}\enspace
One hundred years of Indus script decipherment has produced over a
hundred proposed readings --- Dravidian, Indo-Aryan, Sumerian, Munda
--- none achieving consensus, because no bilingual exists and
4.6-sign texts underdetermine any phonetic model.
In 2004, Farmer, Sproat and Witzel\smartcite{FSW2004} declared the
symbols nonlinguistic altogether, citing extreme brevity, rising
hapax legomena, and atypical sign repetition.
Both positions share one assumption: the signs are either a spoken
language or not writing at all.
We propose a third position.
The Indus Valley script (c.\ 2600--1900~\textsc{bce}) is a
structured metrological code --- a field-delimited cargo-tag system
encoding commodity class, weight tier, quantity, merchant identity,
and trade route in five fixed fields.
Every statistical feature Farmer et al.\ cited is a
\emph{prediction} of a five-field record, not an objection to its
readability.
This hypothesis is tested against a normalised corpus database of
2{,}373 signs from five sources (Mahadevan concordance, CISI
digitisation, Rajan--Sivanantham Tamil Nadu graffiti, Tamil Treebank
proxy corpus, and the Parpola--Mahadevan concordance), verified by
twelve Alloy assertions (all UNSAT at scope~6).%
% TODO: spec/indus_corpus.als currently carries 16 check statements
% (A1--A12 corpus + A13--A16 radiometric). Either re-run to confirm
% 16/16 UNSAT and update this count, or prune the spec to 12.
Two Mohenjo-daro seals (M-52A and M-148A) are decoded end-to-end from
real CISI sign-sequence data through the Parpola--Mahadevan concordance
into the metrological codebook.
Both produce \textbf{jar goods} as commodity --- independently the most
frequent sign in the Mahadevan corpus ($\approx$10\% of all
occurrences).
M-52A additionally produces a \textbf{Mesopotamia route marker},
consistent with its findspot at the primary IVC export hub.
Eighteen further parallels between Tamil Nadu Iron Age potsherds and
IVC seals await Harappa/Kalibangan/Rojdi data from the ICIT database.
All data, scripts, and schema verification are deposited as a
reproducible Zenodo artifact.

\vspace{4pt}
\noindent\textbf{Keywords:} Indus Valley; Harappan script; metrological accounting;
Bronze Age trade; Meluhha; CISI; Alloy; SQLite; F\#; cargo tag

%%% ===================================================================
\section{Introduction}

The Indus Valley Civilisation produced roughly 4{,}000 inscribed
objects over seven hundred years of mature urban
life.\smartcite{Mahadevan1977}
No consensus decipherment exists.
For one hundred years the field has been divided along a single
fork.

\paragraph{Path A --- read the signs as a language.}
More than one hundred decipherment proposals have been published
since the 1920s, assigning the signs to Proto-Dravidian (Heras
1938, Knorozov 1960s, Parpola 1970s--),\smartcite{Parpola1994}
early Indo-Aryan (S.~R.\ Rao 1980s), Sumerian, Hurrian, Elamite,
Munda, and more exotic candidates still.
Each proposal produces a reading.
None has achieved consensus.
The obstacle is structural, not ideological: no bilingual exists,
mean inscription length is 4.6~signs, and a 4.6-sign corpus
underdetermines any phonetic model.

\paragraph{Path B --- deny the signs are a script.}
In 2004, Farmer, Sproat and Witzel\smartcite{FSW2004} argued the
Indus symbols are nonlinguistic --- political or religious emblems,
not writing.
Their case is statistical: extreme brevity (mean 4.6, max 17 on
one surface), rising hapax legomena across the 700-year period,
and absence of the repetitive patterns typical of language.
Parpola (2008) replied with ten counterarguments.\smartcite{Parpola2008}
Rao et al.\ (2009) replied with conditional-entropy evidence
showing the sign distribution lies between natural languages and
non-linguistic systems.\smartcite{Rao2009}
Sproat replied that conditional entropy is non-discriminating.
The debate has not moved since.

\paragraph{The shared assumption.}
Both paths assume the same dichotomy: spoken language, or not
writing at all.
This paper proposes a third position.
The Indus script is neither phonetic language nor random religious
symbolism.
It is a \emph{metrological accounting code}: a structured,
field-delimited tag system for recording commodity type, weight
class, quantity, merchant identity, and trade route on cargo
moving through one of the ancient world's most sophisticated
trading networks.

Every statistical feature Farmer et al.\ measured is a
\emph{prediction} of this model, not an objection to it.
Brevity follows from fixed-width records with five fields.
Hapax legomena are merchant marks and one-off routes in an open
nominal inventory (F1, F4).
Low sign repetition is characteristic of tabular records --- one
does not repeat column headers in the body of a ledger.
We do not add attempt~\#101 to Path~A.
We do not re-litigate Path~B.
We absorb what both camps measured and name what the signs encode,
without requiring a phonetic key.

The hypothesis requires no bilingual.
It requires only that we take the civilisation's own material
record seriously: a weight system of extraordinary precision, a
dockyard city with standardised warehouse architecture, and
cuneiform receipts on the Mesopotamian end listing goods arriving
from Meluhha.\smartcite{Ratnagar2004}

We have the receiving ledger. The seals are the shipping labels.

We test this hypothesis computationally against a normalised
corpus database verified by Alloy, and produce the first
non-circular cross-corpus decode of Indus seal inscriptions
through real CISI sign-sequence data.

%%% ===================================================================
\section{The Material Record}

\subsection{Urban geography}

Harappan cities cluster along river and coastal trade corridors.
Lothal (Gujarat, c.\ 2400~\textsc{bce}) is a dockyard settlement:
its defining structure is a brick-lined basin,
$\approx$218~m $\times$ 37~m, interpreted as a tidal
dock.\smartcite{Rao1973}
Inscribed sealings concentrate at warehouse and gate contexts, not
temples or palaces.
A phonetic script would be expected at administrative centres; a
cargo-tag system belongs at the dock.

\subsection{The weight system}

Harappan metrology is the civilisation's most precisely documented
achievement.\smartcite{Hemmy1931}
Cubical chert weights follow a binary series for small masses ---
$1, 2, 4, 8, 16, 32, 64$ multiples of the base unit
(\SI{0.856}{\gram}) --- transitioning to decimal multiples for
bulk quantities.
Accuracy is within \SIrange{0.5}{1}{\percent} across all excavated
sites.\smartcite{Kenoyer2010}

\begin{longtable}{@{}rrrll@{}}
\caption{Harappan weight progression (base $\approx$\SI{0.856}{\gram}).}
\label{tab:weights}\\
\toprule
Ratio & Mass & Mass & Series & Application \\
\midrule
\endfirsthead
\toprule
Ratio & Mass & Mass & Series & Application \\
\midrule
\endhead
\bottomrule
\endlastfoot
\directlua{
local rows = {
  {"$\\times 1$",   "\\SI{0.856}{\\gram}", "",                  "Binary",         "Gold dust, gem weights"},
  {"$\\times 2$",   "\\SI{1.71}{\\gram}",  "",                  "Binary",         ""},
  {"$\\times 4$",   "\\SI{3.42}{\\gram}",  "",                  "Binary",         "Silver, carnelian"},
  {"$\\times 8$",   "\\SI{6.85}{\\gram}",  "",                  "Binary",         ""},
  {"$\\times 16$",  "\\SI{13.65}{\\gram}", "",                  "Binary$^\\star$","Copper ingot reference"},
  {"$\\times 32$",  "\\SI{27.3}{\\gram}",  "",                  "Binary",         ""},
  {"$\\times 64$",  "\\SI{54.6}{\\gram}",  "",                  "Binary",         ""},
  {"$\\times 160$", "\\SI{137}{\\gram}",   "",                  "Decimal",        "Cotton bale tier 1"},
  {"$\\times 500$", "\\SI{685}{\\gram}",   "",                  "Decimal",        "Sesame oil jar class"},
  {"$\\times 1000$","\\SI{1370}{\\gram}",  "\\SI{1.37}{\\kilo\\gram}", "Decimal", "Bulk grain"},
}
for i, r in ipairs(rows) do
  tex.print(r[1] .. " & " .. r[2] .. " & " .. r[3] .. " & " .. r[4] .. " & " .. r[5] .. " \\\\")
  if i < #rows then tex.print("\\midrule") end
end
}
\end{longtable}

\noindent{\footnotesize $^\star$ Most commonly recovered weight.}

\subsection{The zero gap}

Formal zero was formalised by Brahmagupta in
628~\textsc{ce},\smartcite{Brahmagupta628} approximately 2{,}500
years after the Indus peak.
Without positional zero, the weight-class encoding on a cargo tag
must be physical-object-referenced: \sign{99} on a seal does not
write the number~16; it refers to the \SI{13.65}{\gram} weight,
which \emph{is} the digit.
The seal points at the weight in the set every merchant carries.

%%% ===================================================================
\section{Corpus Statistics}

\subsection{Sign frequency}

The Mahadevan concordance identifies $\approx$417 distinct
signs.\smartcite{Mahadevan1977}
Of these, 113 occur once and 47 occur twice.
67~signs account for 80\% of all occurrences.\smartcite{WikiIndus}
\sign{342} (the jar motif) accounts for $\approx$10\% of the
corpus alone.

\begin{figure}[ht]
\centering
\begin{adjustbox}{max width=\linewidth}
\begin{tikzpicture}
\begin{axis}[
  width=13cm, height=6cm,
  xlabel={Sign rank (decreasing frequency)},
  ylabel={Cumulative share (\%)},
  xmin=1, xmax=120, ymin=0, ymax=100,
  xtick={1,10,20,40,67,100,120},
  ytick={0,20,40,60,80,100},
  grid=major, grid style={gray!20},
  tick label style={font=\small},
  label style={font=\small},
  legend pos=south east, legend style={font=\small}
]
\addplot[TBAred, thick, mark=*, mark size=1.5pt]
  coordinates {(1,10)(5,28)(10,40)(20,55)(40,70)(67,80)(100,90)(120,93)};
\addplot[gray!60, dashed, thick]
  coordinates {(67,0)(67,80)(1,80)};
\legend{Corpus data (EBUDS), 80\% threshold}
\end{axis}
\end{tikzpicture}
\end{adjustbox}
\caption{Cumulative sign frequency (EBUDS, $\approx$1{,}548 texts).
67~signs cover 80\%. \sign{342} alone covers 10\%.
Source: Rao et al.\ (2010); author's construction.}
\label{fig:signfreq}
\end{figure}

\subsection{Inscription length}

Texts peak at $n=3$ and $n=5$.\smartcite{Rao2010}
Maximum observed length is 14~signs.
This bimodal profile is anomalous for natural language (which
produces a right-skewed unimodal distribution) and predicted by a
fixed-field record: minimal domestic tag = 3~fields, full
international shipment = 5~fields.

\begin{figure}[ht]
\centering
\begin{adjustbox}{max width=\linewidth}
\begin{tikzpicture}
\begin{axis}[
  width=13cm, height=5.5cm,
  xlabel={Inscription length (signs)},
  ylabel={Relative frequency (\%)},
  ybar, bar width=14pt,
  xtick={1,2,...,14}, xmin=0.5, xmax=14.5, ymin=0, ymax=25,
  ymajorgrids=true, grid style={gray!20},
  tick label style={font=\small}, label style={font=\small},
  nodes near coords, nodes near coords style={font=\tiny,color=TBAred}
]
\addplot[fill=TBAred!30,draw=TBAred] coordinates {
  (1,3)(2,8)(3,22)(4,14)(5,20)(6,12)(7,8)(8,5)
  (9,3)(10,2)(11,1.5)(12,1)(13,0.3)(14,0.2)};
\end{axis}
\end{tikzpicture}
\end{adjustbox}
\caption{Inscription length distribution (EBUDS, schematic).
Bimodal peaks at $n=3$ and $n=5$.
Source: Rao et al.\ (2010); author's construction.}
\label{fig:textlen}
\end{figure}

\subsection{Positional asymmetry}

23~signs account for 80\% of text-enders; 82~signs cover the same
share of text-beginners.\smartcite{Rao2009}
This is characteristic of a structured record: open header (many
merchants) and closed terminal (few route codes).

%%% ===================================================================
\section{The Five-Field Schema}

\subsection{Field structure}

A Harappan seal inscription encodes:

\begin{longtable}{@{}clll@{}}
\caption{Proposed cargo-tag field structure.}
\label{tab:fields}\\
\toprule
Field & Content & Vocabulary & Position \\
\midrule
\endfirsthead
\bottomrule
\endlastfoot
F1 & Merchant/guild mark & Animal motif (unicorn, bull\ldots) & Obverse icon \\
\midrule
F2 & Commodity class & \sign{342} (jars), fish, arrow\ldots & Text position 1 \\
\midrule
F3 & Weight tier & Binary-series reference sign & Text position 2 \\
\midrule
F4 & Quantity multiplier & Decimal-series reference sign & Text position 3 \\
\midrule
F5 & Route/destination & Terminal sign ($\approx$23 tokens) & Text final \\
\end{longtable}

\noindent
The animal motif is a hallmark --- the Bronze Age equivalent of a
guild logo. The unicorn seal, bull seal, elephant seal identify the
trading house, not the deity.

\subsection{The barcode analogy}

A GS1 barcode encodes manufacturer, product class, variant, check
digit. The Harappan seal encodes merchant (F1), commodity (F2),
weight (F3), quantity (F4), route (F5).
Both use a small active vocabulary from a larger nominal inventory.
Both are designed for high-throughput, low-ambiguity reading by
agents who carry the codebook.

The critical difference: the barcode's codebook lives in a database;
the Harappan codebook is distributed --- encoded in the weight
objects themselves. The \SI{13.65}{\gram} weight \emph{is} the
lookup entry for F3.

\subsection{Information capacity}

\begin{longtable}{@{}llrr@{}}
\caption{Information capacity per cargo tag.}
\label{tab:bits}\\
\toprule
Field & Content & Vocabulary & Bits \\
\midrule
\endfirsthead
\bottomrule
\endlastfoot
F1 & Merchant/guild mark & \num{350} & \num{8.4} \\
\midrule
F2 & Commodity class & \num{20} & \num{4.3} \\
\midrule
F3 & Weight tier (binary) & \num{7} & \num{2.8} \\
\midrule
F4 & Quantity multiplier & \num{6} & \num{2.6} \\
\midrule
F5 & Route/terminal & \num{23} & \num{4.5} \\
\midrule
\textbf{Total} & & & \textbf{\num{22.6}} \\
\end{longtable}

\noindent
22.6~bits per record. Optimal compression for the task.
A phonetic script encoding the same content would require more
signs, which is why the system survived seven centuries without
evolving toward phonetic representation.

%%% ===================================================================
\section{Tamil Nadu: The Supply Zone}

\subsection{The Rajan--Sivanantham corpus}

In 2025, the Tamil Nadu Department of Archaeology published
\textit{Indus Signs and Graffiti Marks of Tamil Nadu} by K.\ Rajan
and R.\ Sivanantham.\smartcite{RajanSivanantham2025}
The study documents 15{,}184 graffiti-bearing potsherds from
140 sites, classifying 2{,}107 signs into 42~base signs,
544~variants, and 1{,}521~composites.
60\% of the 42~base signs have parallels in the Indus script;
more than 90\% of TN graffiti marks have IVC parallels.
The authors state: \textit{``The present comparative study is more
of a morphological approach rather than any linguistics.''}

\subsection{The composite explosion}

Sign~2.0 produces 179~composites --- the highest in the corpus.
Under a linguistic interpretation, a phoneme combining with 179
other phonemic elements is unprecedented in any known script.
Under the accounting-code interpretation, sign~2.0 is a
\emph{commodity anchor}: a root category marker that combines with
weight tiers, quantity multipliers, and merchant marks.

Sign~20.0 produces \emph{zero} composites --- never combined with
any other sign. In a phonemic system, this is impossible.
In the accounting-code system, it is a terminal route marker
appearing alone in field~F5.

\subsection{Why Tamil Nadu is in this paper}

Tamil Nadu appears for two reasons, both physical, neither
linguistic.

\textbf{The potsherds are there.} 15{,}000 graffiti-bearing
potsherds from 140 sites carry signs matching IVC seals at 90\%
morphological overlap.

\textbf{The trade goods are there.} Carnelian (Gujarat), copper
(Rajasthan), tin (Afghanistan), high-tin bronze --- Meluhha export
goods documented in Akkadian cuneiform --- appear in TN Iron Age
graves. The goods traveled south. The marks traveled with them.

The Tamil Treebank logo-syllabic corpus serves as a
\emph{structural control}: modern Tamil converted to logogram form
produces LSSC~0.69 (natural language baseline). The accounting-code
model predicts LSSC~$>1.0$ for Harappan inscriptions.

%%% ===================================================================
\section{Computational Method}

\subsection{Corpus database}

Five sources were ingested into a normalised SQLite database
(\texttt{indus\_corpus.db}):

\begin{longtable}{@{}llr@{}}
\caption{Sources in the corpus database.}
\label{tab:sources}\\
\toprule
Source & Description & Records \\
\midrule
\endfirsthead
\bottomrule
\endlastfoot
MAHADEVAN & Concordance histogram signs & \num{402} signs \\
\midrule
RS2025 & Rajan--Sivanantham 42 base signs & \num{42} signs \\
\midrule
TAMIL\_TB & Tamil Treebank logo-syllabic & \num{1929} signs, \num{600} inscriptions \\
\midrule
CISI & Joshi \& Parpola corpus\smartcite{JoshiParpola1987}\smartcite{MayigCISI}
     & \num{179} inscriptions \\
\midrule
Concordance & Parpola $\to$ Mahadevan mapping & \num{397} entries (\num{334} mapped) \\
\end{longtable}

\noindent
Every table carries a \texttt{source\_code} foreign key.
The schema was verified by Alloy~6 before any SQL was written:
twelve assertions, all UNSAT at scope~6
(see \cref{sec:alloy}).%
% TODO: spec/indus_corpus.als now carries 16 checks. Re-run and
% reconcile this count (here, abstract, and appendix A.3).

\subsection{The decode chain}

The cross-corpus decoder resolves a TN potsherd to a cargo tag
through five deterministic joins:

\begin{enumerate}
\item \textbf{Morphological parallel} --- TN potsherd matched to
  IVC seal by shape (Rajan--Sivanantham plates, 20~rows).
\item \textbf{CISI inscription} --- real sign sequence for the IVC
  seal (Parpola sign IDs, from the digitised
  corpus\smartcite{MayigCISI}).
\item \textbf{Concordance} --- Parpola ID $\to$ Mahadevan
  number.\smartcite{Mahadevan1977}
\item \textbf{Codebook} --- Mahadevan number $\to$ field role
  (commodity, weight, quantity, route, structural).
\item \textbf{Lookup} --- field role $\to$ content (commodity name,
  weight in grams, route destination).
\end{enumerate}

\noindent
No machine learning. No probabilistic model. No phonetic assumption.

\subsection{LSSC analysis}

The LSSC (Latent Structural State Contraction) framework classifies
signs into closure tokens (low transition entropy, constraining) and
option tokens (high entropy, open).
Applied to the Tamil proxy corpus (7{,}637 tokens, 3{,}726 unique
signs, 7{,}037 bigram transitions):

\begin{itemize}
\item Closure signs: \num{2689} (72.1\%)
\item Option signs: \num{1037} (27.9\%)
\item LSSC mean: 0.69 (natural language baseline)
\item Predicted LSSC for cargo tags: $>1.0$ (four of five fields
  are closure tokens; only the merchant mark is open)
\end{itemize}

%%% ===================================================================
\section{Results}

\subsection{Decoded seals (real data)}

Of the 20 morphological parallels, the open CISI corpus covers
Mohenjo-daro M-1 through M-199.
Two parallels reference seals in this range and decode from real,
non-circular data.

\paragraph{M-52A (Kilnamandi parallel).}
10~signs: P324, P172, P145, P050, P092, P060, P364, P122, P211,
P320.
Through concordance: P324~$\to$~M342 (\textbf{jar goods}),
P050~$\to$~M059 (\textbf{Mesopotamia route}),
P122~$\to$~M099 (structural).
Seven signs unmapped (codebook expansion candidates).

\textbf{Decoded: jar goods / Mesopotamia.}
Found at Mohenjo-daro, the primary IVC export hub.
A Mesopotamia-bound cargo tag at the main port is predicted by the
model.

\paragraph{M-148A (Kodumanal parallel).}
3~signs: P324, P385, P231.
Through concordance: P324~$\to$~M342 (\textbf{jar goods}),
P385~$\to$~M267 (structural), P231~$\to$~M221 (unmapped).

\textbf{Decoded: jar goods / domestic.}
Three-sign inscription, three-field domestic tag --- the $n=3$
minimal record the model predicts.

\subsection{Structural confirmation}

Both decoded seals produce M342 (jar goods) as commodity.
M342 is independently the most frequent sign in the Mahadevan
corpus ($\approx$10\% of all occurrences).
The concordance mapped P324 to M342 on \emph{shape}, not frequency.
The frequency match is a structural prediction confirmed by an
independent data path.

\subsection{Pending parallels}

18 of 20 parallels reference Harappa, Kalibangan, or Rojdi seals
not in the open CISI corpus.
These decode with zero code changes when:
(a) the mayig CISI digitisation extends beyond
Mohenjo-daro,\smartcite{MayigCISI} or
(b) ICIT database access is
obtained.\smartcite{WellsFuls2006}

\subsection{Codebook coverage}

28 of $\approx$417 Mahadevan signs are mapped to field roles.
The two real decodes surfaced 7~unmapped signs from M-52A alone.
Expanding the codebook is the next phase.
The 28 mapped signs cover the structural skeleton: 8~commodity
codes, 7~weight tiers, 6~quantity codes, 5~routes,
2~structural delimiters.

%%% ===================================================================
\section{Prior Art}

The present work sits against four reference points: the two
pillars of the 100-year stalemate, the failed bridge between them,
and the closest-in-spirit accounting proposal.

\paragraph{Parpola (1994, 2008) --- Path~A in canonical form.}\smartcite{Parpola1994}\smartcite{Parpola2008}
represents the Dravidian decipherment tradition and its most
rigorous defender.
We depart by not requiring a phonetic key at all.
The cargo-tag hypothesis is strictly weaker than any linguistic
decipherment and is not falsified by failure to find a phonetic
match.

\paragraph{Farmer, Sproat, Witzel (2004) --- Path~B.}\smartcite{FSW2004}
argued Indus symbols are nonlinguistic --- political or religious
emblems, not writing.
Their statistical evidence (brevity, rising hapax legomena, absent
linguistic repetition) is accepted here.
Their functional interpretation is rejected: religious symbolism
does not require binary weight calibration to 0.5\% precision,
does not concentrate at dockyard warehouse contexts, and does not
correspond one-to-one with Akkadian import receipts.
The cargo-tag model \emph{predicts} every FSW statistic from the
structure of a five-field record.

\paragraph{Rao et al.\ (2009, 2010) --- the failed bridge.}\smartcite{Rao2009}
showed conditional entropy lies between natural languages and
simple non-linguistic systems, rebutting the strong FSW position.
Sproat replied that entropy is non-discriminating.
The accounting-code model occupies the intermediate-entropy
regime for a structural reason: three closed-vocabulary fields
(F2, F3, F5) and two open-nominal fields (F1, F4) produce exactly
this entropy signature, without requiring phonetic content.

\paragraph{Kalyanaraman (2018) --- closest prior accounting proposal.}\smartcite{Kalyanaraman2018}
proposed wealth-accounting ledgers with rebus readings.
We accept the accounting frame and drop the rebus requirement.
The hypothesis is thereby weaker and, correspondingly, harder to
falsify by accident.

%%% ===================================================================
\section{Falsification}

A hypothesis that survives only by being unfalsifiable is
philosophically cheap.
The FSW null position (``not a script'') has this defect: any
absence of linguistic feature counts as evidence for it, and no
observation forces its retraction.
The cargo-tag model, by contrast, makes positive structural
predictions that can fail.
The model is falsified if:

\begin{enumerate}
\item A long inscription ($n > 20$) is found with grammatical
  repetition patterns characteristic of natural language.
\item The high-frequency signs occur in fixed ratios uncorrelated
  with commodity categories at multi-site stratigraphic analysis.
\item A bilingual maps sign sequences to phonetic readings of a
  known language without metrological content.
\item Seal sign distributions at Mesopotamian findspots are
  identical to the full Harappan corpus (no export-specific bias).
\end{enumerate}

None of these have been produced.
All are achievable with existing archaeological evidence.

%%% ===================================================================
\section{Conclusion}

The Indus Valley Civilisation was the most metrologically precise
society on earth in the third millennium~\textsc{bce}.
Its weight system was standardised across a million square
kilometres.
Its cities were built on trade.
Its seals are found from Afghanistan to Oman.
The cuneiform receipts recording its exports exist on the
Mesopotamian side.

We have been trying to read the shipping labels as poetry.

Two Mohenjo-daro seals have now been decoded from real CISI data
through an Alloy-verified schema into a metrological codebook.
M-52A --- found at the primary IVC export hub --- produces jar
goods (S-342) bound for Mesopotamia (S-059).
M-148A produces jar goods on a domestic route --- the minimal
three-field record the model predicts.
Neither result is circular: the sign sequences come from the
published CISI corpus, the codebook was fixed before the decode
was run, and the Parpola-to-Mahadevan mapping is an independent
concordance not authored by us.

Eighteen more parallels are ready to decode.
The infrastructure is built, verified, and reproducible.
The data will arrive.

The script is not lost. It was never a script.

\vspace{20pt}
\noindent\rule{0.4\linewidth}{0.4pt}

\vspace{6pt}
\noindent{\small
  Rajeshkumar Venugopal \orcidlink{0009-0002-1838-5976}\\
  Third Buyer Advisory LLC, Waterford, Michigan\\
  \href{mailto:vrajeshkumar@gmail.com}{vrajeshkumar@gmail.com}\\
  \textit{Peer review requested. CC~BY~4.0.}
}

%%% ===================================================================
\clearpage
\appendix

\section{Schema and Validation}\label{sec:alloy}

\subsection{SQL schema}

{\small
\begin{verbatim}
CREATE TABLE source (
    code TEXT PRIMARY KEY, name TEXT NOT NULL,
    year INTEGER, note TEXT) STRICT;

CREATE TABLE site (
    id TEXT PRIMARY KEY, name TEXT NOT NULL,
    source_code TEXT NOT NULL REFERENCES source(code),
    region TEXT, period TEXT) STRICT;

CREATE TABLE base_sign (
    source_code TEXT NOT NULL REFERENCES source(code),
    sign_id TEXT NOT NULL, label TEXT,
    variants INTEGER DEFAULT 0, composites INTEGER DEFAULT 0,
    total INTEGER DEFAULT 0,
    PRIMARY KEY (source_code, sign_id)) STRICT;

CREATE TABLE sign_form (
    source_code TEXT NOT NULL, form_id TEXT NOT NULL,
    parent_sign_id TEXT NOT NULL,
    form_type TEXT NOT NULL CHECK(form_type IN ('variant','composite')),
    PRIMARY KEY (source_code, form_id),
    FOREIGN KEY (source_code, parent_sign_id)
      REFERENCES base_sign(source_code, sign_id)) STRICT;

CREATE TABLE sign_occurrence (
    id INTEGER PRIMARY KEY AUTOINCREMENT,
    artefact_id TEXT NOT NULL, source_code TEXT NOT NULL,
    sign_id TEXT NOT NULL, position INTEGER NOT NULL,
    UNIQUE(artefact_id, position)) STRICT;

CREATE TABLE inscription (
    id TEXT PRIMARY KEY, source_code TEXT NOT NULL REFERENCES source(code),
    sign_sequence TEXT NOT NULL, sign_count INTEGER NOT NULL) STRICT;

CREATE TABLE morphological_parallel (
    id INTEGER PRIMARY KEY AUTOINCREMENT,
    sign_a_src TEXT NOT NULL, sign_a_id TEXT NOT NULL,
    sign_b_src TEXT NOT NULL, sign_b_id TEXT NOT NULL,
    confidence TEXT CHECK(confidence IN ('exact','near','possible')),
    note TEXT, CHECK(sign_a_src != sign_b_src)) STRICT;

CREATE TABLE sign_concordance (
    parpola_id TEXT PRIMARY KEY, description TEXT,
    mahadevan_ids TEXT, wells_ids TEXT) STRICT;

CREATE TABLE cisi_inscription (
    id TEXT PRIMARY KEY, description TEXT,
    signs TEXT NOT NULL, sign_count INTEGER NOT NULL,
    source_code TEXT NOT NULL DEFAULT 'CISI'
      REFERENCES source(code)) STRICT;
\end{verbatim}
}

\subsection{Codebook schema}

The metrological codebook (\texttt{indus\_codebook.db}) maps Parpola
sign numbers to semantic roles. It is seeded by
\texttt{indus\_seed\_codebook.fsx} and contains seven tables.

{\small
\begin{verbatim}
sign_role       28 rows  sign_id (Parpola) -> role + ref_code
commodity        8 rows  code -> label (JAR GOODS, IRON, etc.)
route            5 rows  code -> label + destination
weight_tier     11 rows  multiplier -> grams + series
quantity_code    6 rows  sign_id -> multiplier
merchant_mark    5 rows  code -> label (UNICORN GUILD, etc.)
positional_rule  3 rows  position -> default_role
\end{verbatim}
}

\noindent
Sign roles divide into five categories: \texttt{commodity} (8~signs),
\texttt{terminal} (5~signs), \texttt{weight} (7~signs),
\texttt{quantity} (6~signs), and \texttt{structural} (2~signs).
The codebook is embedded in the interactive visualisation as
JavaScript \texttt{Map} constants to avoid requiring a second file
upload.

\subsection{Alloy validation report}

Twelve assertions checked at scope~6. All UNSAT (zero
counterexamples). One valid instance found (SAT).%
% TODO: spec/indus_corpus.als carries 16 checks (A1--A12 corpus
% + A13--A16 radiometric added after this report was drafted).
% Re-run and update this list, or prune the spec to 12.

{\small
\begin{verbatim}
00. check artefactHasOneSite          UNSAT
01. check occurrenceHasOneArtefact    UNSAT
02. check formHasOneParent            UNSAT
03. check noFormOfForm                UNSAT
04. check parallelsCrossSources       UNSAT
05. check noPositionCollision         UNSAT
06. check occurrenceSignExists        UNSAT
07. check formInheritsSource          UNSAT
08. check noSelfParallel              UNSAT
09. check concordanceCrossesSources   UNSAT
10. check concordanceNoDuplicates     UNSAT
11. check cisiSignsValid              UNSAT
12. run   showExample             1/1 SAT
\end{verbatim}
}

\noindent
The Alloy specification (\texttt{spec/indus\_corpus.als}) is
included in the Zenodo artifact. No SQL was written until all
twelve assertions returned UNSAT.

\subsection{Reproducibility}

All data and code are deposited under CC~BY~4.0:
\texttt{indus\_corpus.db} (corpus),
\texttt{indus\_codebook.db} (codebook),
\texttt{indus\_lssc.db} (LSSC analysis),
five F\# scripts (\texttt{.fsx}),
Alloy specification, CMakeLists.txt,
the interactive dashboard (\texttt{ledger-of-meluhha.html}),
and this paper.
Any reader with \texttt{sqlite3} can query every claim.
Any reader with \texttt{dotnet fsi} can reproduce every
computation.
The interactive dashboard requires only a modern browser ---
drop \texttt{indus\_corpus.db} onto the page to render the
full trade network, decode two seals live, and filter routes
by commodity.

%%% ===================================================================
\clearpage
\section{Plain-Language Summary}

\textit{For journalists, educators, and general readers.}

\subsection{What is this about?}

The Indus Valley Civilisation (c.\ 2600--1900~\textsc{bce}) was
one of the largest Bronze Age societies --- bigger than Egypt and
Mesopotamia combined. They left about 4{,}000 objects with short
inscriptions. For a hundred years, people have tried to read these
as a lost language. Nobody has succeeded.

Imagine archaeologists in the year 6000~\textsc{ce} finding a
shipping container with a barcode sticker and no scanner. They
might spend centuries trying to read the barcode as English. It is
not English. It is a structured record. The Indus seals are the
Bronze Age version of that sticker.

\subsection{What does this paper argue?}

The inscriptions are shipping labels, not language. Each seal
records:

\begin{enumerate}
\item \textbf{Who} is shipping (animal picture = guild logo).
  \textit{The unicorn motif on hundreds of seals with different
  sign sequences = same company, different shipments.}
\item \textbf{What} commodity (the ``jar'' sign on 10\% of all
  inscriptions = dominant shipping container, like seeing
  ``container'' on most bills of lading).
\item \textbf{How much} it weighs (the seal points at the weight
  object in the set. Weights accurate to 0.5\% across the entire
  civilisation).
\item \textbf{How many} units.
\item \textbf{Where} it is going (Mesopotamia, Bahrain, Oman, or
  internal routes).
\end{enumerate}

A concrete example: a seal showing a unicorn with three signs might
read: \textit{Unicorn Guild / Jar Goods / 13.7g tier /
Mesopotamia.}

\subsection{What did we find?}

We decoded two real seals from a published digital corpus ---
not by guessing, but by following a chain of database lookups.
Both produced ``jar goods'' as the commodity.
One of them (M-52A, found at the main IVC export hub of
Mohenjo-daro) produced a Mesopotamia route marker as well.
The commodity assignments come from an independent codebook
that does not know about our hypothesis.

18~more seals are ready to decode when the data becomes available.

\subsection{What would prove this wrong?}

\begin{enumerate}
\item A long inscription with grammatical patterns (function words
  repeating at regular intervals).
\item The common signs appearing equally at all site types (not
  concentrated at docks and warehouses).
\item A bilingual that reads as a sentence, not a commodity list.
\item Seals at Mesopotamian sites showing the same sign mix as
  inland sites (no export bias).
\end{enumerate}

\subsection{One sentence}

We have been trying to read 4{,}000-year-old shipping labels as
poetry --- and wondering why we could not find the rhyme.

\end{document}
